\documentclass[10pt,a4paper]{amsart}
\usepackage{amsmath,amssymb,amsthm,graphicx}
% Retained preamble for the Stacks Project chapter {CH} (shared project preamble).
% The chapter file pinned-original-source.tex only carries \input{{preamble}} on its first line,
% so no preamble text lives inside the chapter itself.  The Stacks Project uses one single shared
% preamble for every chapter: the snapshot stages it three times (sources/stacks-chow/preamble.tex,
% sources/stacks-etale-cohomology/preamble.tex and focused-seeds/AGP/sources/stacks-snapshot-20261011/
% preamble.tex) and those copies are identical in size.  The block below is copied mechanically,
% line for line, from sources/stacks-chow/preamble.tex lines 14-18 and 154-217: the Xy-pic setup and
% the native theorem declarations together with the mathematical macros.  Nothing was added, removed
% or redefined; the wrapper class and the amsmath/amssymb/amsthm/graphicx base packages are supplied
% by the assembler.
% For commutative diagrams we use Xy-pic
% Chapter: stacks-cohomology.tex -- Cohomology of Sheaves.
\usepackage[all]{xy}

% We use 2cell for 2-commutative diagrams.
\xyoption{2cell}

% Theorem environments.
%
\theoremstyle{plain}
\newtheorem{theorem}[subsection]{Theorem}
\newtheorem{proposition}[subsection]{Proposition}
\newtheorem{lemma}[subsection]{Lemma}

\theoremstyle{definition}
\newtheorem{definition}[subsection]{Definition}
\newtheorem{example}[subsection]{Example}
\newtheorem{exercise}[subsection]{Exercise}
\newtheorem{situation}[subsection]{Situation}

\theoremstyle{remark}
\newtheorem{remark}[subsection]{Remark}
\newtheorem{remarks}[subsection]{Remarks}

\numberwithin{equation}{subsection}

% Macros
%
\def\lim{\mathop{\mathrm{lim}}\nolimits}
\def\colim{\mathop{\mathrm{colim}}\nolimits}
\def\Spec{\mathop{\mathrm{Spec}}}
\def\Hom{\mathop{\mathrm{Hom}}\nolimits}
\def\Ext{\mathop{\mathrm{Ext}}\nolimits}
\def\SheafHom{\mathop{\mathcal{H}\!\mathit{om}}\nolimits}
\def\SheafExt{\mathop{\mathcal{E}\!\mathit{xt}}\nolimits}
\def\Sch{\mathit{Sch}}
\def\Mor{\mathop{\mathrm{Mor}}\nolimits}
\def\Ob{\mathop{\mathrm{Ob}}\nolimits}
\def\Sh{\mathop{\mathit{Sh}}\nolimits}
\def\NL{\mathop{N\!L}\nolimits}
\def\CH{\mathop{\mathrm{CH}}\nolimits}
\def\proetale{{pro\text{-}\acute{e}tale}}
\def\etale{{\acute{e}tale}}
\def\QCoh{\mathit{QCoh}}
\def\Ker{\mathop{\mathrm{Ker}}}
\def\Im{\mathop{\mathrm{Im}}}
\def\Coker{\mathop{\mathrm{Coker}}}
\def\Coim{\mathop{\mathrm{Coim}}}

% Boxtimes
%
\DeclareMathSymbol{\boxtimes}{\mathbin}{AMSa}{"02}

%
% Macros for moduli stacks/spaces
%
\def\QCohstack{\mathcal{QC}\!\mathit{oh}}
\def\Cohstack{\mathcal{C}\!\mathit{oh}}
\def\Spacesstack{\mathcal{S}\!\mathit{paces}}
\def\Quotfunctor{\mathrm{Quot}}
\def\Hilbfunctor{\mathrm{Hilb}}
\def\Curvesstack{\mathcal{C}\!\mathit{urves}}
\def\Polarizedstack{\mathcal{P}\!\mathit{olarized}}
\def\Complexesstack{\mathcal{C}\!\mathit{omplexes}}
% \Pic is the operator that assigns to X its picard group, usage \Pic(X)
% \Picardstack_{X/B} denotes the Picard stack of X over B
% \Picardfunctor_{X/B} denotes the Picard functor of X over B
\def\Pic{\mathop{\mathrm{Pic}}\nolimits}
\def\Picardstack{\mathcal{P}\!\mathit{ic}}
\def\Picardfunctor{\mathrm{Pic}}
\def\Deformationcategory{\mathcal{D}\!\mathit{ef}}
\usepackage{cleveref}
\crefname{originallemma}{original lemma}{original lemmas}
\Crefname{originallemma}{Original lemma}{Original lemmas}
\makeatletter
\newlabel{derived-lemma-derived-canonical-delta-functor}{{0152}{source}{}{sourcepointer-1}{}}
\newlabel{derived-lemma-derived-canonical-delta-functor@cref}{{[originallemma][0][]0152}{[1][0][]source}{}{}{}}
\makeatother
\crefname{originalsection}{original section}{original sections}
\Crefname{originalsection}{Original section}{Original sections}
\makeatletter
\newlabel{derived-section-canonical-delta-functor}{{014Z}{source}{}{sourcepointer-2}{}}
\newlabel{derived-section-canonical-delta-functor@cref}{{[originalsection][0][]014Z}{[1][0][]source}{}{}{}}
\makeatother
\crefname{originaltheorem}{original theorem}{original theorems}
\Crefname{originaltheorem}{Original theorem}{Original theorems}
\makeatletter
\newlabel{injectives-theorem-K-injective-embedding-grothendieck}{{079P}{source}{}{sourcepointer-3}{}}
\newlabel{injectives-theorem-K-injective-embedding-grothendieck@cref}{{[originaltheorem][0][]079P}{[1][0][]source}{}{}{}}
\makeatother
\crefname{originallemma}{original lemma}{original lemmas}
\Crefname{originallemma}{Original lemma}{Original lemmas}
\makeatletter
\newlabel{lemma-mayer-vietoris}{{01EB}{source}{}{sourcepointer-4}{}}
\newlabel{lemma-mayer-vietoris@cref}{{[originallemma][0][]01EB}{[1][0][]source}{}{}{}}
\makeatother
\crefname{originallemma}{original lemma}{original lemmas}
\Crefname{originallemma}{Original lemma}{Original lemmas}
\makeatletter
\newlabel{lemma-restrict-K-injective-to-open}{{08BS}{source}{}{sourcepointer-5}{}}
\newlabel{lemma-restrict-K-injective-to-open@cref}{{[originallemma][0][]08BS}{[1][0][]source}{}{}{}}
\makeatother
\title{Mayer-Vietoris distinguished triangle for the derived pushforward along the inclusions of two open subspaces}
\author{The Stacks Project Authors}
\date{stacks-project:cohomology.tex:local-snapshot-20261011; GNU FDL 1.2 or later}
\begin{document}
\maketitle
\noindent\textit{Editorial note.} All delivered mathematical text is the author's own, copied from the pinned Stacks Project chapter file cohomology.tex; no statement or proof step was rewritten, reordered or removed.\par\medskip
\noindent\textit{Editorial note.} The proof defers the exactness of the sequence of sections to the proof of the Mayer-Vietoris lemma, which is a navigation pointer; the termwise statement that makes the appeal applicable is written out inside the counted span.\par\medskip
\noindent\textit{Editorial note.} The counted spans carry 5 navigation stubs and no bibliography citation.\par\medskip
\noindent\textit{Editorial note.} The preamble is the chapter's retained minimal preamble (sources/stacks-cohomology/retained-preamble.tex lines 1-81), a verbatim excerpt of the native preamble.\par\medskip
\noindent\textit{Editorial note.} The label lemma-exact-sequence-j-star is defined only in cohomology.tex, so no cross-chapter restatement had to be weighed for this card.\par\medskip
\noindent\textit{Allowed standard prerequisites (editorial).}\begin{itemize}
\item K-injective resolutions with injective terms exist. Injectives, Theorem injectives-theorem-K-injective-embedding-grothendieck, invoked at proof lines 8136-8137. Supplies the representative by which the derived pushforwards are computed.
\item Restriction to an open subspace preserves K-injectivity. Same chapter, Lemma lemma-restrict-K-injective-to-open, invoked at proof lines 8138-8139. Supplies that the restricted complexes compute the derived pushforwards on U, V and the intersection.
\item Distinguished triangle attached to a short exact sequence of complexes. Derived Categories, Section derived-section-canonical-delta-functor and Lemma derived-lemma-derived-canonical-delta-functor, invoked at proof lines 8141-8145. Supplies the triangle once the short exact sequence of complexes is in hand.
\item The termwise exact sequence of sections over W, W cap U, W cap V and W cap U cap V. Same chapter, Lemma lemma-mayer-vietoris, whose proof is cited for this exactness at proof lines 8154-8163. Supplies the exactness of the short exact sequence of complexes of sections.
\end{itemize}\medskip
\ifdefined\hypertarget\hypertarget{sourcepointer-1}{}\fi
\noindent\textit{Source reference.} Original lemma 0152: sources/stacks-chow/tags.txt, source lines 1493--1493. Pointer outside the counted span, invoked once for context; the load-bearing argument of the counted statement is the author's proof reproduced in full in the proof span.\par\medskip
\ifdefined\hypertarget\hypertarget{sourcepointer-2}{}\fi
\noindent\textit{Source reference.} Original section 014Z: sources/stacks-chow/tags.txt, source lines 1490--1490. Pointer outside the counted span, invoked once for context; the load-bearing argument of the counted statement is the author's proof reproduced in full in the proof span.\par\medskip
\ifdefined\hypertarget\hypertarget{sourcepointer-3}{}\fi
\noindent\textit{Source reference.} Original theorem 079P: sources/stacks-chow/tags.txt, source lines 9279--9279. Pointer outside the counted span, invoked once for context; the load-bearing argument of the counted statement is the author's proof reproduced in full in the proof span.\par\medskip
\ifdefined\hypertarget\hypertarget{sourcepointer-4}{}\fi
\noindent\textit{Source reference.} Original lemma 01EB: sources/stacks-cohomology/pinned-original-source.tex, source lines 767--767. Pointer outside the counted span, invoked once for context; the load-bearing argument of the counted statement is the author's proof reproduced in full in the proof span.\par\medskip
\ifdefined\hypertarget\hypertarget{sourcepointer-5}{}\fi
\noindent\textit{Source reference.} Original lemma 08BS: sources/stacks-cohomology/pinned-original-source.tex, source lines 7840--7840. Pointer outside the counted span, invoked once for context; the load-bearing argument of the counted statement is the author's proof reproduced in full in the proof span.\par\medskip
\begin{lemma}
\label{lemma-exact-sequence-j-star}
Let $(X, \mathcal{O}_X)$ be a ringed space.
Let $X = U \cup V$ be the union of two open subspaces.
For any object $E$ of $D(\mathcal{O}_X)$ we have a distinguished
triangle
$$
E \to 
Rj_{U, *}E|_U \oplus Rj_{V, *}E|_V \to
Rj_{U \cap V, *}E|_{U \cap V} \to
E[1]
$$
in $D(\mathcal{O}_X)$.
\end{lemma}

\begin{proof}
Choose a K-injective complex $\mathcal{I}^\bullet$ representing $E$
whose terms $\mathcal{I}^n$ are injective objects of
$\textit{Mod}(\mathcal{O}_X)$, see Injectives, Theorem
\ref{injectives-theorem-K-injective-embedding-grothendieck}.
We have seen that $\mathcal{I}^\bullet|U$ is a K-injective complex
as well (Lemma \ref{lemma-restrict-K-injective-to-open}). Hence
$Rj_{U, *}E|_U$ is represented by $j_{U, *}\mathcal{I}^\bullet|_U$.
Similarly for $V$ and $U \cap V$. Hence the distinguished triangle
of the lemma is the distinguished triangle associated (by
Derived Categories, Section
\ref{derived-section-canonical-delta-functor} and especially
Lemma \ref{derived-lemma-derived-canonical-delta-functor})
to the short exact sequence of complexes
$$
0 \to
\mathcal{I}^\bullet \to
j_{U, *}\mathcal{I}^\bullet|_U \oplus j_{V, *}\mathcal{I}^\bullet|_V \to
j_{U \cap V, *}\mathcal{I}^\bullet|_{U \cap V} \to
0.
$$
This sequence is exact because for any $W \subset X$ open
and any $n$ the sequence
$$
0 \to
\mathcal{I}^n(W) \to
\mathcal{I}^n(W \cap U) \oplus \mathcal{I}^n(W \cap V) \to
\mathcal{I}^n(W \cap U \cap V) \to
0
$$
is exact (see proof of Lemma \ref{lemma-mayer-vietoris}).
\end{proof}

\end{document}
